% Options for packages loaded elsewhere
% Options for packages loaded elsewhere
\PassOptionsToPackage{unicode}{hyperref}
\PassOptionsToPackage{hyphens}{url}
\PassOptionsToPackage{dvipsnames,svgnames,x11names}{xcolor}
%
\documentclass[
  english,
  11pt,
]{article}
\usepackage{xcolor}
\usepackage[margin=25mm]{geometry}
\usepackage{amsmath,amssymb}
\setcounter{secnumdepth}{-\maxdimen} % remove section numbering
\usepackage{iftex}
\ifPDFTeX
  \usepackage[T1]{fontenc}
  \usepackage[utf8]{inputenc}
  \usepackage{textcomp} % provide euro and other symbols
\else % if luatex or xetex
  \usepackage{unicode-math} % this also loads fontspec
  \defaultfontfeatures{Scale=MatchLowercase}
  \defaultfontfeatures[\rmfamily]{Ligatures=TeX,Scale=1}
\fi
\usepackage{lmodern}
\ifPDFTeX\else
  % xetex/luatex font selection
\fi
% Use upquote if available, for straight quotes in verbatim environments
\IfFileExists{upquote.sty}{\usepackage{upquote}}{}
\IfFileExists{microtype.sty}{% use microtype if available
  \usepackage[]{microtype}
  \UseMicrotypeSet[protrusion]{basicmath} % disable protrusion for tt fonts
}{}
\makeatletter
\@ifundefined{KOMAClassName}{% if non-KOMA class
  \IfFileExists{parskip.sty}{%
    \usepackage{parskip}
  }{% else
    \setlength{\parindent}{0pt}
    \setlength{\parskip}{6pt plus 2pt minus 1pt}}
}{% if KOMA class
  \KOMAoptions{parskip=half}}
\makeatother
% Make \paragraph and \subparagraph free-standing
\makeatletter
\ifx\paragraph\undefined\else
  \let\oldparagraph\paragraph
  \renewcommand{\paragraph}{
    \@ifstar
      \xxxParagraphStar
      \xxxParagraphNoStar
  }
  \newcommand{\xxxParagraphStar}[1]{\oldparagraph*{#1}\mbox{}}
  \newcommand{\xxxParagraphNoStar}[1]{\oldparagraph{#1}\mbox{}}
\fi
\ifx\subparagraph\undefined\else
  \let\oldsubparagraph\subparagraph
  \renewcommand{\subparagraph}{
    \@ifstar
      \xxxSubParagraphStar
      \xxxSubParagraphNoStar
  }
  \newcommand{\xxxSubParagraphStar}[1]{\oldsubparagraph*{#1}\mbox{}}
  \newcommand{\xxxSubParagraphNoStar}[1]{\oldsubparagraph{#1}\mbox{}}
\fi
\makeatother


\usepackage{longtable,booktabs,array}
\newcounter{none} % for unnumbered tables
\usepackage{calc} % for calculating minipage widths
% Correct order of tables after \paragraph or \subparagraph
\usepackage{etoolbox}
\makeatletter
\patchcmd\longtable{\par}{\if@noskipsec\mbox{}\fi\par}{}{}
\makeatother
% Allow footnotes in longtable head/foot
\IfFileExists{footnotehyper.sty}{\usepackage{footnotehyper}}{\usepackage{footnote}}
\makesavenoteenv{longtable}
\usepackage{graphicx}
\makeatletter
\newsavebox\pandoc@box
\newcommand*\pandocbounded[1]{% scales image to fit in text height/width
  \sbox\pandoc@box{#1}%
  \Gscale@div\@tempa{\textheight}{\dimexpr\ht\pandoc@box+\dp\pandoc@box\relax}%
  \Gscale@div\@tempb{\linewidth}{\wd\pandoc@box}%
  \ifdim\@tempb\p@<\@tempa\p@\let\@tempa\@tempb\fi% select the smaller of both
  \ifdim\@tempa\p@<\p@\scalebox{\@tempa}{\usebox\pandoc@box}%
  \else\usebox{\pandoc@box}%
  \fi%
}
% Set default figure placement to htbp
\def\fps@figure{htbp}
\makeatother



\ifLuaTeX
\usepackage[bidi=basic,shorthands=off]{babel}
\else
\usepackage[bidi=default,shorthands=off]{babel}
\fi
\ifLuaTeX
  \usepackage{selnolig} % disable illegal ligatures
\fi


\setlength{\emergencystretch}{3em} % prevent overfull lines

\providecommand{\tightlist}{%
  \setlength{\itemsep}{0pt}\setlength{\parskip}{0pt}}



 


\makeatletter
\@ifpackageloaded{caption}{}{\usepackage{caption}}
\AtBeginDocument{%
\ifdefined\contentsname
  \renewcommand*\contentsname{Table of contents}
\else
  \newcommand\contentsname{Table of contents}
\fi
\ifdefined\listfigurename
  \renewcommand*\listfigurename{List of Figures}
\else
  \newcommand\listfigurename{List of Figures}
\fi
\ifdefined\listtablename
  \renewcommand*\listtablename{List of Tables}
\else
  \newcommand\listtablename{List of Tables}
\fi
\ifdefined\figurename
  \renewcommand*\figurename{Figure}
\else
  \newcommand\figurename{Figure}
\fi
\ifdefined\tablename
  \renewcommand*\tablename{Table}
\else
  \newcommand\tablename{Table}
\fi
}
\@ifpackageloaded{float}{}{\usepackage{float}}
\floatstyle{ruled}
\@ifundefined{c@chapter}{\newfloat{codelisting}{h}{lop}}{\newfloat{codelisting}{h}{lop}[chapter]}
\floatname{codelisting}{Listing}
\newcommand*\listoflistings{\listof{codelisting}{List of Listings}}
\makeatother
\makeatletter
\makeatother
\makeatletter
\@ifpackageloaded{caption}{}{\usepackage{caption}}
\@ifpackageloaded{subcaption}{}{\usepackage{subcaption}}
\makeatother
\usepackage{bookmark}
\IfFileExists{xurl.sty}{\usepackage{xurl}}{} % add URL line breaks if available
\urlstyle{same}
% fallback for those not using the hyperref driver hyperxmp:
\makeatletter
\@ifundefined{xmpquote}{\newcommand{\xmpquote}[1]{#1}}{}
\makeatother
\hypersetup{
  pdftitle={The vortex radius is not the only small length},
  pdflang={en},
  colorlinks=true,
  linkcolor={blue},
  filecolor={Maroon},
  citecolor={Blue},
  urlcolor={Blue},
  pdfcreator={LaTeX via pandoc}}


\title{The vortex radius is not the only small length}
\usepackage{etoolbox}
\makeatletter
\providecommand{\subtitle}[1]{% add subtitle to \maketitle
  \apptocmd{\@title}{\par {\large #1 \par}}{}{}
}
\makeatother
\subtitle{Why fine ripples and molecular motion need their own rulers}
\author{}
\date{2026-09-12}
\begin{document}
\maketitle


\href{../paper-map/article.html}{Return to the paper map} · Physical
interpretation, article 2 of 3

Imagine a painting that is one metre wide but contains lines a tenth of
a millimetre apart. The canvas size does not tell you how fine a camera
must be to resolve the lines.

A vortex has the same distinction. Its core radius measures the size of
an intense region. Fine waves inside that region can vary over a much
shorter distance.

This matters when asking whether a continuum fluid model has a physical
interpretation. We need a ruler for the variation we want to describe,
not just a ruler for the whole object.

\subsection{Compare the right lengths}\label{compare-the-right-lengths}

A one-millimetre core might, in an illustrative scenario, contain
changes over one micrometre. That is a thousand times smaller.

The overall radius may look comfortably large while the fine structure
demands much smaller averaging regions.

The paper has several lengths: the core width, the axial extent, the
width of a pulse envelope, and the spacing of the oscillation inside it.
A further distinction is the distance over which velocity changes
appreciably, which also depends on amplitude and direction.

We cannot replace all of these by ``the vortex size.'' Nor have our
physical scale studies reconstructed the complete pulse wavelengths and
gradients.

\subsection{An averaging box must be neither too large nor too
small}\label{an-averaging-box-must-be-neither-too-large-nor-too-small}

To assign a local velocity, imagine averaging molecular motion in a
small box.

If the box is too large, it blends different parts of the flow and
erases the feature we wanted to describe. If it is too small, it
contains too few molecules for the intended smooth average.

Suppose a box contains 1,000 molecules on average. Shrink each side by
ten while holding the density fixed. Its volume is divided by 1,000, so
its expected molecular count drops to one.

That is why a modest-looking change in the required linear resolution
can cause a large change in molecular statistics.

\begin{figure}

\centering{

\pandocbounded{\includegraphics[keepaspectratio,alt={A larger averaging cube labelled 1000 molecules is compared with a cube one tenth as wide labelled one molecule.}]{figures/reader.png}}

}

\caption{\label{fig-average}At fixed density, reducing each side of an
averaging cube by ten reduces its expected molecular count by a
thousand. The count illustrates volume scaling, not a universal minimum
number for fluid modelling.}

\end{figure}%

The whole core may still contain many molecules. The question concerns
the small regions needed to describe its local variations.

\subsection{Gas molecules bring another
ruler}\label{gas-molecules-bring-another-ruler}

Between collisions, gas molecules travel a typical distance called the
\textbf{mean free path}.

If the flow changes over a distance much larger than that path, a local
continuum description has a useful separation of scales. If the two
become comparable, molecules can cross significantly different
conditions between collisions.

A common preliminary comparison is

\[
\text{length ratio}=
\frac{\text{mean free path}}
{\text{flow-variation length}}.
\]

This is the Knudsen number. For the stated room-temperature air
reference, the mean free path is about \textbf{67 nanometres}. A
variation length about one hundred times larger is approximately
\textbf{6.7 micrometres}.

The corresponding ratio of 0.01 is often used as an engineering screen.
It is not an exact universal boundary or proof that every continuum
assumption holds below it.

\subsection{Speed and length ask different physical
questions}\label{speed-and-length-ask-different-physical-questions}

Comparing speed with sound speed asks whether density changes may
matter. Comparing molecular flight length with flow-variation length
asks whether a local continuum description is appropriate.

Either concern can arise first, depending on the flow. If a fine ripple
is much smaller than the core radius, the molecular-length screen can be
reached earlier than a calculation using the core radius would suggest.

A gas can also remain a continuum while needing a \textbf{compressible}
model. ``Continuum'' and ``incompressible'' are different assumptions,
so they should not be treated as a single switch.

For a liquid, the gas mean-free-path picture is not directly
transferable. Molecular spacing, correlations, relaxation, and
thermodynamic state require their own treatment.

\subsection{Molecules also need time to
respond}\label{molecules-also-need-time-to-respond}

Lengths are only part of the story. The local state may change so
quickly that molecules do not have enough time to establish the kind of
local balance assumed by the model.

A rough flow time is ``variation length divided by speed.'' Shorten the
distance or increase the speed and that time decreases.

Comparing it with a molecular response time is another useful question.
Our gas calculations use a declared collision-time proxy. They do not
establish a universal liquid response time or a physical realization of
the paper's detailed field.

\subsection{What our regime map actually
compares}\label{what-our-regime-map-actually-compares}

The supporting study follows several assumed variation lengths along the
same illustrative radius--speed curve. It asks how the order of the
physical screening markers changes.

These are scenarios. The figure does not measure the actual smallest
length in the paper's pulses. Its purpose is to expose the dependency
that a single core-radius plot hides.

The physical discussion draws on the scalings in Sections 2--3 of the
paper and the property/model references recorded in the
\href{source-notes.md}{source notes}.

\subsection{Where this fits in the bigger
picture}\label{where-this-fits-in-the-bigger-picture}

\textbf{A large object can contain structure too fine for the model used
to describe it.} Physical interpretation therefore needs both the
overall geometry and the actual variation scales.

This is a companion question about the meaning of the continuum
construction, not another proof block. Continue with
\href{../d10_compressible_vortex/article.html}{what changes when gas
density is allowed to vary}, or return to the
\href{../paper-map/article.html}{paper map}.

The \href{technical-article.html}{technical article},
\href{../../code/d9_regime_map/study.ipynb}{notebook}, and
\href{source-notes.md}{source notes} retain the scenario assumptions,
units, and unresolved gradient scales.




\end{document}
